% ==============================================================================
\section{Related Work}
\label{sec:related}
% ==============================================================================

The AHPP lies at the intersection of four research streams that we briefly
review here: aircraft hangar maintenance scheduling with spatial layout,
broader aircraft maintenance scheduling at the fleet or line level,
airport ground operations and stand or gate allocation, and methodological
precedents from MILP scheduling with time-window or blocking constraints
and from GRASP metaheuristics for combined assignment and scheduling
problems.  The first stream is the closest in application domain and
supplies our main point of comparison; the second positions the AHPP
within aircraft maintenance optimisation at large; the third provides
adjacent airport assignment and recovery models; the fourth contributes
the modelling and algorithmic toolkit.  We do not address long-term
aircraft maintenance routing or check scheduling, which decide when
aircraft enter maintenance over a planning horizon; the AHPP starts from
a set of accepted maintenance requests and optimises their
spatial-temporal placement inside a single hangar.

The closest body of work is the line of MILP formulations for the
integrated hangar maintenance scheduling and layout problem.
\citet{Qin2017} introduce the first MILP that couples maintenance
scheduling and multi-period parking layout with non-overlapping
constraints between aircraft.  Building on that line, \citet{Qin2018}
develop a two-stage exact-heuristic procedure for aircraft parking
stand allocation that trades off expected revenue against safety
distance, and \citet{Qin2018b} extend it with revised No-Fit Polygons
and a family of heuristic-based valid inequalities for the safety-margin
objective.  \citet{Qin2019} broaden the view to joint maintenance
scheduling and parking under a hangar with time-varying capacity and
blocking along taxi routes, proposing an event-based discrete-time MILP
and a rolling-horizon heuristic for large instances, and
\citet{Qin2020} integrate scheduling, layout, movement trajectories and
multi-skill technician assignment under an MRO outsourcing mode through
a two-stage decomposition.  More recently, the preprint of
\citet{Pazhooh2025} proposes a continuous-time MILP for integrated
hangar scheduling and layout, reporting substantial computational
improvements over previous event-based formulations and experiments on
larger instances.  These works share with the AHPP the
assign-and-schedule decision structure but model the spatial
interactions through geometric primitives --- No-Fit Polygons for
non-overlapping placement, taxi-route capacities, safety distances, or
overlap of continuous-time windows --- so that the operational impact
of moving an aircraft is captured indirectly via route feasibility,
parking-plan feasibility, or movement-path constraints, rather than as
an explicit cost term induced by an asymmetric, directed relation
between positions.

Stepping out of the hangar but staying in aircraft maintenance,
\citet{Shaukat2020} formulate the line maintenance scheduling problem
as a MIP that places maintenance work in the layovers between scheduled
flights through integrated exact and sequential heuristic approaches;
\citet{Witteman2021} cast multi-year task allocation across an
airline's fleet as a time-constrained variable-sized bin-packing
problem and solve 45-aircraft instances with a worst-fit-decreasing
heuristic; and \citet{TorresSanchez2020} provide a multi-objective
MILP that integrates airline fleet maintenance scheduling with tail
assignment and workshop capacity over 30-day horizons with 529
aircraft and eight workshops.  These works tackle rich temporal and
resource structures at the fleet or network scale but abstract away
the within-hangar geometry that drives the AHPP: positions are pooled
into workshop capacities or maintenance slots, not laid out as a graph
of directed blocking relationships.  The AHPP is therefore
complementary, operating at the operational time scale of a single
hangar.

At the broader scale of airport ground operations, the systematic
review of \citet{Dahanayaka2026} surveys 199 articles on gate and stand
allocation, turnaround scheduling, workforce assignment, and surface
movement.  Among the relevant models, \citet{Guepet2015} solve the
airport stand allocation problem with both an exact MILP and spatial
and time decomposition heuristics; their setting is the closest
methodological cousin to the AHPP --- aircraft-to-position assignment
plus timing, with shadow restrictions and tow movements as soft
objective terms --- and their two major European airport datasets
illustrate the same exact-versus-decomposition trade-off that drives
our solver portfolio.  \citet{Asadi2021} formulate the gate-assignment
schedule-recovery problem with constrained resources;
\citet{Evler2021a, Evler2021b} extend the approach to recover airline
ground operations under uncertain arrival times; and \citet{Yao2024}
couple stand allocation with ground-support-vehicle scheduling.  All
these works share with the AHPP the aircraft-to-position assignment
layer and the soft-time-window machinery, but operate at the apron
level with shared servicing resources rather than on a hangar with
directed blocking between maintenance positions.

The AHPP also draws on a broader literature on MILP scheduling with
time windows and blocking-like interactions outside the airport
domain.  \citet{Mao2018} address multi-degree and multi-hoist cyclic
scheduling with time windows, where hoists share a track and must
avoid spatial collisions; their disjunctive-time-window approach is
structurally similar to our blocking-arc constraints.
\citet{Lee2024} model multi-satellite mission and communication
scheduling as a MILP with time-window and conflict constraints, and
\citet{Baya2022} apply MILP to a grain facility with shared storage
capacity.  None of these works expresses interference as a directed
asymmetric relation tied to a position-assignment decision, but they
illustrate the modelling techniques --- time windows, conflict graphs,
big-$M$ disjunctions --- that the AHPP combines.  On the algorithmic
side, the topology-aware heuristic developed in this paper is a Greedy
Randomized Adaptive Search Procedure; its foundations go back to
\citet{Feo1995} and are surveyed in the annotated bibliographies of
\citet{FestaResende2002, FestaResende2009}.  Closer to our setting,
\citet{Wang2024} attack a yard-crane scheduling problem with dynamic
assignment of input/output points using construction and local-search
operators that mirror our coupling between position selection and
start-time choice, and \citet{Ferone2018} provide the methodological
basis for the biased-randomisation step in our $\alpha$-controlled
restricted candidate list.

Against this background, the present paper contributes along two
fronts.  On the \emph{modelling} side, it abstracts the hangar layout
as a directed blocking-arc graph on maintenance positions, making the
asymmetry between front and rear --- and its dependence on the
assignment decision --- an explicit feature of the problem rather than
an indirect consequence of taxi routes, safety distances, or
geometric overlap as captured by No-Fit Polygons
\citep{Qin2018, Qin2018b, Qin2019, Qin2020, Pazhooh2025}.  The entry
and exit manoeuvres induced by this graph then become an optimisable
objective component through a compact six-binary indicator structure,
so that movements are minimised jointly with delay and makespan
through user-defined weights, rather than appearing only as
feasibility side-effects of capacities or distances.

On the \emph{algorithmic} side, the paper combines an aircraft-level
MILP --- where each aircraft is a single uninterrupted maintenance
block --- with a topology-aware GRASP and a fixed-assignment scheduler
glued together by a safe pipeline that returns the best solution found
by its components.  The exact solver is competitive on small
instances, in line with the scale at which the formulations of
\citet{Guepet2015, Shaukat2020, Pazhooh2025} also close to optimality,
while the safe pipeline takes over as the fleet size grows: the
crossover between exact and heuristic methods in our benchmark sits
near $R = 14$ aircraft for a 60-second wall-clock budget, consistent
with the $14$--$25$-aircraft threshold reported across the hangar
literature for exact MILP optimality \citep{Pazhooh2025, Guepet2015,
Shaukat2020, Qin2019}.  Rather than competing directly with
general continuous-time layout formulations on their own benchmark
classes, the focus is on a topology-driven blocking abstraction and on
a fast solver portfolio that delivers high-quality feasible solutions
on a controlled benchmark of congested hangar layouts, in the regime
where the exact solver leaves large optimality gaps.
